\begin{abstract}
Mobile GUI agents now handle multi-step tasks on familiar apps, but generalizing
to applications never seen during training remains the bottleneck: recent work
shows that online reinforcement learning improves a vision-language agent on
unseen task instances by large margins yet barely moves the needle on unseen
\emph{applications}, and points to few-shot test-time adaptation as a promising
but preliminary remedy. We develop that direction into \textbf{EvoFSM-RL}, a
test-time adaptation framework that, when deploying to a new app on a small
adaptation budget, co-adapts two channels in a single loop: it evolves a
structured in-context prior \emph{and} updates the policy weights. The prior is a
two-layer finite state machine whose lower layer is app-specific and whose upper
layer is a category-generic, linter-enforced unit of cross-app transfer; the loop
mutates the upper layer via a frozen reference language model and updates a LoRA
adapter via group-relative policy optimization, and the same loop runs at
source-pool pretraining and at target-app deployment. We evaluate at two levels
of generalization that share this method. \emph{Within} a benchmark
(AndroidWorld+), under a protocol that is simultaneously pool-, category-, and
template-disjoint, static category knowledge lifts near-transfer success by
$+9.3$~pp over a zero-shot baseline and online upper-layer evolution adds another
$+3.7$~pp, with a far-transfer panel acting as a built-in null control.
\emph{Across} benchmarks (AndroidWorld+ to an independently authored MobileWorld
deployment target), we find that priors which transfer within a benchmark
collapse across the benchmark boundary, and trace the collapse to interpretable
causes that motivate adapting the prior on the target rather than transplanting
it. The paper contributes a method for test-time adaptation to unseen apps, a
joint prompt-and-weight adaptation algorithm, and a multi-level benchmark
partition for measuring mobile-agent generalization.
\end{abstract}
